\chapter{Conclusion} \label{chap:ch6}

In the following sections, we will summarize the contributions of this dissertation
and discuss the Research Questions that were addressed. 

\section{Contributions} \label{sec:se61}

Throughout this dissertation, we have presented \ac{DisRUPT}, a model agnostic testbed 
that enables the evaluation of distributed inference for \acp{MLP} and \acp{CNN} 
models given a set of devices and a configuration file containing the model 
distribution, established routes and input data.

We have also highlighted a set of strategies and good practices for the deployment
of \ac{DNN} models in the edge environment, that lead to a lower inference latency 
and network energy consumption. Namely, the impact that communication steps have on
the distributed inference process, and the importance of minimizing their occurrence,
whether by reducing the number of partitions or by selecting certain partition
points that minimize the amount of data that needs to be transmitted.

Taking into account the initial goal and research questions, as stated in Chapter~\ref{chap:intro},
we can conclude that this work addressed all the research questions, as follows:

\begin{itemize}
    \item \textbf{Hypothesis: } \textit{"Is it possible to create a reliable tool to 
        support the deployment of complex \ac{DNN} models on far-edge and edge devices?"}
        
        \subitem This hypothesis was confirmed, as the \ac{DisRUPT} supports exactly that,
        even if limited in supported model architectures and width, by the underlying 
        NAMP compatibility and the available resources of the devices.  
    
    \item \textbf{RQ1: } \textit{"Is it feasible to have collaborative ML inferences 
        using edge and far-edge devices without compromising performance?"}

        \subitem This research question was denied, although our work
        shows that the mentioned process is feasible, the performance is compromised
        in the sense that the inference latency is increased, even if not significantly. 

    \item \textbf{RQ2: } \textit{"What is the most efficient method to enable reliable 
        cooperation between edge and far-edge devices partaking in an inference 
        network?"}

        \subitem During the experiments, we found out that, from the cooperative
        distribution strategies, there was no method that was generally more efficient
        than the others. And that efficiency is highly dependent on the model 
        architecture, network conditions, and the available resources. 

    \item \textbf{RQ3: } \textit{"How efficiently can the resources of edge and 
        far-edge devices be utilized towards collaborative inference tasks?"}

        \subitem The experiments showed that the resources of edge and far-edge
        devices can be efficiently utilized, although the intermediate software layer
        that is required to enable inference on \acp{MCU} highly limits the available
        resources, and therefore, the efficiency of the process.
        
\end{itemize}


\section{Future Work} \label{sec:se62}

Even though the contributions of this dissertation reached the expected goals, the work
can be further extended in several ways that further evaluate the feasibility of
having \ac{DNN} execution in the edge environment and in which forms. 

Firstly, the compatibility of the \ac{DisRUPT} testbed can be expanded to support
more \ac{DNN} architectures, such as recurrent and attention layers, to
enable the partitioning of more complex and capable models. Likewise, the
support of intra-layer partitioning can be added to the testbed, removing the
current limitation of only supporting inter-layer partitioning, and model width
constraints.

Additionally, the testbed can be extended to support devices with accelerators, such 
as \acp{TPU}, \acp{NPU}, and \acp{GPU}, allowing the evaluation of a more diverse 
range of strategies and deployment scenarios, to further explore the potential of
distributed inference in the edge environment. 

Finally, the orchestration and scheduling strategies can be further improved by
implementing a real time computing scheduler that can dynamically adapt to the
network conditions and the available resources. This results in quality improvements
on the collected data, which will be closer to a real-world scenario, and therefore, 
more reliable and useful for future research in this field.